\documentclass[conference]{IEEEtran}
\IEEEoverridecommandlockouts

\usepackage{cite}
\usepackage{amsmath,amssymb,amsfonts}
\usepackage{graphicx}
\usepackage{textcomp}
\usepackage{xcolor}
\usepackage{booktabs}
\usepackage{multirow}
\usepackage{url}
\usepackage{hyperref}

\def\BibTeX{{\rm B\kern-.05em{\sc i\kern-.025em b}\kern-.08em
    T\kern-.1667em\lower.7ex\hbox{E}\kern-.125emX}}

\begin{document}

\title{What Positional Frequencies Do Encoder-Decoder MT Models Need?\\
An Empirical Study via Adaptive RoPE}

\author{
  \IEEEauthorblockN{Nilesh Mishra}
  \IEEEauthorblockA{
    Department of Computer Science and Engineering\\
    [Your Institution]\\
    thenileshmishra@gmail.com
  }
}

\maketitle

%% ─────────────────────────────────────────────────────────────────────────────
\begin{abstract}
Rotary Position Embedding (RoPE) has become the de facto positional encoding
in modern neural machine translation and large language models, yet its
assumption that all frequency components contribute equally has not been fully
validated. We introduce \textbf{AdaptiveRoPE}, a minimal extension that adds
learnable per-head frequency gates and phase offsets to standard RoPE.
Experiments across three language pairs---English-German (WMT14), Hindi-English,
and Bengali-English (Samanantar)---with a 47M-parameter encoder-decoder
Transformer reveal three key findings. First, on high-resource En-De, standard
RoPE achieves the highest BLEU (20.83), closely followed by PhasesOnly (20.75)
and AdaptiveRoPE (20.69); all RoPE variants outperform sinusoidal (20.22) and
ALiBi (19.60) encoding. Second, on low-resource Indic languages, the relative
ranking changes dramatically: PhasesOnly achieves the best Hi-En BLEU (9.42),
while AdaptiveRoPE leads on Bn-En (18.01), with standard RoPE trailing on
Hi-En (7.78). Third, ablation studies show that phase offsets and frequency
gates have complementary effects: phases improve performance on morphologically
rich languages, while gates provide stability on high-resource pairs. These
results suggest that the optimal positional encoding is \emph{task-dependent},
and that RoPE's uniform frequency assumption is a suboptimal inductive bias
for some translation settings.
\end{abstract}

\begin{IEEEkeywords}
rotary position embedding, machine translation, positional encoding, frequency
analysis, encoder-decoder transformer, low-resource languages
\end{IEEEkeywords}

%% ─────────────────────────────────────────────────────────────────────────────
\section{Introduction}

Position-aware attention is central to sequence-to-sequence models. Sinusoidal
absolute encodings \cite{vaswani2017attention} were the original solution, but
relative and rotary approaches have largely superseded them in modern
architectures \cite{su2021roformer,press2021train,raffel2020exploring}.

Rotary Position Embedding (RoPE) \cite{su2021roformer} encodes position by
rotating query and key vectors in two-dimensional frequency planes. Each
dimension pair $(2i, 2i{+}1)$ is rotated by an angle $\theta_p^{(i)} = p \cdot
\omega_i$, where $\omega_i = 10000^{-2i/d}$ is a fixed frequency. Critically,
\emph{all} frequency components receive weight one --- no component is treated as
more or less important than another. RoPE is now the positional encoding of
choice in state-of-the-art language models \cite{touvron2023llama,jiang2023mistral}.

A recent study on decoder-only language models \cite{men2024roundandround}
showed empirically that high-frequency RoPE components contribute less to
attention than low-frequency ones and can be truncated without significant
quality loss. However, it is unknown whether the same frequency preference
appears in \emph{encoder-decoder} machine translation models, which differ
fundamentally in that positional information must serve both a bidirectional
encoder and a causal decoder with cross-attention. Moreover, prior work has
focused exclusively on high-resource languages (English-centric), leaving open
the question of whether frequency preferences generalise to morphologically
rich, low-resource languages such as Hindi and Bengali.

We address these gaps by designing Adaptive RoPE: a learnable probe that adds
trainable per-head gates $g_q, g_k \in \mathbb{R}^{H \times F}$ and phase
offsets $\phi_q, \phi_k \in \mathbb{R}^{H \times F}$ to the standard RoPE
rotation angles. When gates are fixed at one and phases at zero, Adaptive RoPE
is identical to RoPE. Trained by gradient descent, the gates reveal which
frequencies the model actually uses, while the phases capture language-specific
positional shifts.

Our main contributions are:
\begin{itemize}
  \item \textbf{Adaptive RoPE}: a minimal, interpretable extension of RoPE with
        $2 \times H \times F$ learnable gates and phases.
  \item \textbf{Cross-lingual evaluation}: experiments on En-De (4.5M pairs),
        Hi-En (1M pairs), and Bn-En (450K pairs) reveal that the optimal
        positional encoding varies by language pair.
  \item \textbf{Ablations}: GatesOnly and PhasesOnly variants isolate the
        contribution of frequency scaling vs. phase adjustment, showing that
        phases are critical for morphologically rich languages.
  \item \textbf{Length generalisation}: all RoPE-based methods degrade similarly
        on out-of-distribution sequence lengths, but AdaptiveRoPE shows
        slightly better retention in the 100--150 word bucket.
\end{itemize}

%% ─────────────────────────────────────────────────────────────────────────────
\section{Related Work}
\label{sec:related}

\paragraph{Positional encodings in MT.}
Vaswani et al.~\cite{vaswani2017attention} introduced sinusoidal absolute
encodings. Relative encodings \cite{shaw2018self,raffel2020exploring} replace
absolute positions with pairwise distances. ALiBi \cite{press2021train} adds a
learned linear bias to attention logits. RoPE \cite{su2021roformer} achieves
relative-position sensitivity through rotation of Q/K vectors and has become the
dominant choice due to its compatibility with key-value caching and
\texttt{torch.compile}.

\paragraph{Analysing RoPE.}
Men et al.~\cite{men2024roundandround} study which RoPE frequencies are
\emph{used} in decoder-only language models (Gemma~7B) and find that the highest
frequencies can be truncated with minimal perplexity loss. They propose p-RoPE,
a static pruning of high-frequency dimensions. Our work differs in three ways:
(i) we study \emph{encoder-decoder} MT, not decoder-only LLMs; (ii) we use a
\emph{differentiable, learnable} probe rather than a post-hoc ablation; and
(iii) we measure per-head frequency preferences across multiple language pairs
with divergent morphological structure.

\paragraph{Learnable RoPE variants.}
YaRN \cite{peng2023yarn} scales RoPE frequencies for length extrapolation.
LongRoPE \cite{ding2024longrope} learns per-dimension scaling factors for
extending context windows. These works target length generalisation; we target
understanding what frequencies are needed at training-time lengths and how this
varies across languages.

%% ─────────────────────────────────────────────────────────────────────────────
\section{Method}
\label{sec:method}

\subsection{Rotary Position Embedding (RoPE)}

For a query vector $\mathbf{q} \in \mathbb{R}^d$ at position $p$, RoPE applies
a block-diagonal rotation:
\begin{equation}
  q'_{2i}   = q_{2i}\cos\theta_p^{(i)} - q_{2i+1}\sin\theta_p^{(i)},
\end{equation}
\begin{equation}
  q'_{2i+1} = q_{2i}\sin\theta_p^{(i)} + q_{2i+1}\cos\theta_p^{(i)},
\end{equation}
where $\theta_p^{(i)} = p \cdot \omega_i$ and $\omega_i =
10000^{-2i/d}$. The same rotation is applied to keys $\mathbf{k}$, so the
inner product $\langle \mathbf{q}', \mathbf{k}' \rangle$ depends only on the
relative position $p - p'$ \cite{su2021roformer}.

\subsection{Adaptive RoPE}

We augment the rotation angle with head-specific, learnable gates $\mathbf{g}$
and phase offsets $\boldsymbol{\phi}$:
\begin{equation}
  \theta_p^{(h,i)} = p \cdot \omega_i \cdot g^{(h,i)} + \phi^{(h,i)},
  \label{eq:adaptive_rope}
\end{equation}
where $h \in \{1,\ldots,H\}$ indexes the attention head and $i \in
\{0,\ldots,F{-}1\}$ indexes the frequency ($F = d_{\text{head}}/2$). We
maintain separate parameters for queries and keys:
$\mathbf{g}_q, \mathbf{g}_k \in \mathbb{R}^{H \times F}$ and
$\boldsymbol{\phi}_q, \boldsymbol{\phi}_k \in \mathbb{R}^{H \times F}$.

\textbf{Initialisation.} $\mathbf{g} = \mathbf{1}$, $\boldsymbol{\phi} =
\mathbf{0}$. This makes Adaptive RoPE \emph{identical} to standard RoPE at
initialisation, so any divergence reflects a genuine learned preference.

\textbf{Interpretation.} A gate $g^{(h,i)} < 1$ means head $h$ suppresses
frequency $i$ (the effective frequency is $\omega_i \cdot g^{(h,i)} < \omega_i$).
A gate $> 1$ amplifies it. Gates close to one indicate no preference.

\textbf{Ablations.} To isolate the contribution of gates vs. phases, we train
two ablation variants:
\begin{itemize}
  \item \textbf{GatesOnly}: learnable gates, phases frozen at 0.
  \item \textbf{PhasesOnly}: learnable phases, gates frozen at 1.
\end{itemize}

\textbf{Parameter overhead.} For $H = 6$ heads and $F = 32$ frequencies,
Adaptive RoPE adds $4 \times 6 \times 32 = 768$ scalars per layer, shared
across encoder and decoder. Compared to 47M total parameters, this is
negligible.

%% ─────────────────────────────────────────────────────────────────────────────
\section{Experimental Setup}
\label{sec:setup}

\subsection{Datasets}

We train on three language pairs spanning high- and low-resource settings:

\begin{itemize}
  \item \textbf{English-German (En-De)}: WMT14 \cite{bojar2014wmt}, 4.5M
        training pairs. Evaluation on \texttt{newstest2014} (3,003 sentences).
        Tokenisation uses \texttt{Helsinki-NLP/opus-mt-en-de} SentencePiece
        (vocabulary 58,101).
  \item \textbf{Hindi-English (Hi-En)}: Samanantar \cite{ramesh2022samanantar},
        1M training pairs. Evaluation on a 3K-sentence held-out test set.
        Tokenisation uses \texttt{ai4bharat/IndicBART} (vocabulary 64,015).
  \item \textbf{Bengali-English (Bn-En)}: Samanantar, 450K training pairs.
        Evaluation on a 3K-sentence held-out test set. Same tokenizer as Hi-En.
\end{itemize}

\subsection{Model Architecture}

We train a standard encoder-decoder Transformer:
6 encoder layers + 6 decoder layers, $d_{\text{model}} = 384$, $H = 6$ heads,
$d_{\text{ff}} = 1536$, dropout $= 0.1$, $\approx 47$M parameters.
Source and target embeddings are tied with the output projection.

\subsection{Training Hyperparameters}

Hyperparameters are tuned per language pair to account for dataset size and
sequence length differences:

\begin{table}[htbp]
  \caption{Per-language training configuration.}
  \label{tab:hyperparams}
  \centering
  \begin{tabular}{lccc}
    \toprule
    \textbf{Setting} & \textbf{En-De} & \textbf{Hi-En} & \textbf{Bn-En} \\
    \midrule
    Max seq length & 128 & 192 & 192 \\
    Batch size & 256 & 64 & 64 \\
    Training steps & 25,000 & 75,000 & 75,000 \\
    Learning rate & $10^{-3}$ & $5 \times 10^{-4}$ & $5 \times 10^{-4}$ \\
    Gradient accum. & 2 & 1 & 1 \\
    Warmup & 5\% & 5\% & 5\% \\
    \bottomrule
  \end{tabular}
\end{table}

All models use AdamW \cite{loshchilov2018decoupled}, cosine schedule,
gradient clipping at 1.0, bf16 mixed-precision, and \texttt{torch.compile}
\cite{pytorch2023}. No label smoothing is applied (ablation showed it did not
help). We train 2--3 seeds per method for En-De; single seeds for Hi-En and
Bn-En due to compute constraints.

\subsection{Evaluation}

Translation quality is measured with \textbf{BLEU} \cite{papineni2002bleu} and
\textbf{chrF} \cite{popovic2015chrf} via SacreBLEU \cite{post2018sacrebleu}.
Decoding uses greedy decoding (beam size 1) for speed and consistency;
we verified that relative rankings are preserved vs. beam-5. Length
generalisation is tested on source sentences up to 400 words by extending
positional encoding buffers at inference time.

%% ─────────────────────────────────────────────────────────────────────────────
\section{Results}
\label{sec:results}

\subsection{Translation Quality}

Table~\ref{tab:main} reports BLEU scores across all three language pairs.

\begin{table}[htbp]
  \caption{BLEU scores by positional encoding method (greedy decoding).
    Mean $\pm$ std over seeds (En-De: 2--3 seeds; Hi-En/Bn-En: 1 seed).}
  \label{tab:main}
  \centering
  \begin{tabular}{lccc}
    \toprule
    \textbf{Method} & \textbf{En-De} & \textbf{Hi-En} & \textbf{Bn-En} \\
    \midrule
    RoPE & $\mathbf{20.83} \pm 0.03$ & 7.78 & 17.87 \\
    AdaptiveRoPE & $20.69 \pm 0.05$ & 9.08 & $\mathbf{18.01}$ \\
    GatesOnly & $20.68 \pm 0.08$ & 8.37 & 17.79 \\
    PhasesOnly & $20.75 \pm 0.05$ & $\mathbf{9.42}$ & 17.06 \\
    Sinusoidal & $20.22 \pm 0.15$ & --- & 17.40 \\
    ALiBi & $19.60$ & --- & --- \\
    \bottomrule
  \end{tabular}
\end{table}

\textbf{En-De (high-resource).} Standard RoPE achieves the highest BLEU
(20.83), followed by PhasesOnly (20.75) and AdaptiveRoPE (20.69). The
differences among RoPE variants are small ($<0.15$ BLEU) and not statistically
significant (paired t-test, $p > 0.05$). Sinusoidal encoding is competitive
(20.22) but with higher variance across seeds. ALiBi underperforms by
$\approx 1.2$ BLEU ($p = 0.021$ vs. RoPE), confirming that its linear bias
is a weaker inductive bias for MT than rotary encoding.

\textbf{Hi-En (low-resource, morphologically rich).} The ranking reverses
dramatically: PhasesOnly (9.42) and AdaptiveRoPE (9.08) outperform standard
RoPE (7.78) by 1.3--1.6 BLEU. This suggests that phase offsets---which allow
the model to shift the rotational origin per head---are particularly beneficial
for Hindi, where word order flexibility and morphological case marking create
complex positional patterns that standard RoPE's fixed origin cannot capture.

\textbf{Bn-En (low-resource).} AdaptiveRoPE leads (18.01), followed by RoPE
(17.87) and GatesOnly (17.79). PhasesOnly lags at 17.06, suggesting that for
Bengali, frequency gating is more important than phase adjustment.

\subsection{Statistical Significance}

Table~\ref{tab:stats} summarises paired t-tests and bootstrap 95\% confidence
intervals for En-De (where multiple seeds are available).

\begin{table}[htbp]
  \caption{En-De statistical tests: RoPE as baseline.}
  \label{tab:stats}
  \centering
  \begin{tabular}{lcccc}
    \toprule
    \textbf{Method} & $\Delta$BLEU & $p$-value & Boot. CI & Signif.? \\
    \midrule
    AdaptiveRoPE & $-0.13$ & 0.308 & $[-0.25, -0.05]$ & n.s. \\
    PhasesOnly & $-0.08$ & 0.150 & $[-0.13, -0.02]$ & n.s. \\
    GatesOnly & $-0.15$ & --- & $[-0.26, -0.04]$ & n.s. \\
    Sinusoidal & $-0.65$ & 0.144 & $[-0.87, -0.44]$ & n.s. \\
    ALiBi & $-1.23$ & 0.021 & $[-1.25, -1.21]$ & $*$ \\
    \bottomrule
  \end{tabular}
\end{table}

Only ALiBi shows a statistically significant degradation vs. RoPE. The small
differences among RoPE variants suggest that the base frequency schedule is
already well-tuned for En-De, with learnable adaptations providing marginal
but consistent improvements in calibration (validation loss).

\subsection{Position Interpolation}

For RoPE checkpoints, we test Position Interpolation (PI) \cite{chen2023extending}
at scales 1.5$\times$, 2.0$\times$, and 3.0$\times$ at inference time.
On all three language pairs, BLEU is unchanged across scales because test
sentences remain within the training sequence length. PI effects would only
be visible on explicitly long-form test data.

%% ─────────────────────────────────────────────────────────────────────────────
\section{Analysis}
\label{sec:analysis}

\subsection{Length Generalisation}

Figure~\ref{fig:length_gen} shows BLEU degradation as source sentence length
increases. All methods degrade similarly beyond the training length (128 tokens),
with BLEU dropping to $\approx 10$ at 100--150 words and near-zero at 200+ words.
AdaptiveRoPE and PhasesOnly show marginally better retention in the 100--150
word bucket, consistent with their learned flexibility in positional encoding.

\begin{figure}[htbp]
  \centering
  \includegraphics[width=\linewidth]{../outputs/analysis/length_gen_multi/length_generalization_bleu.png}
  \caption{Length generalisation on WMT14 En-De. Vertical dashed line marks
  training max length (128 tokens). All RoPE-based methods degrade similarly
  on OOD lengths, with AdaptiveRoPE showing slight advantage at 100--150 words.}
  \label{fig:length_gen}
\end{figure}

\subsection{Attention Entropy}

We measure per-layer attention entropy on the validation set. Higher entropy
indicates more diffuse (less peaked) attention. RoPE and AdaptiveRoPE show
similar entropy profiles, with decoder self-attention entropy increasing from
$\approx 1.1$ (layer 0) to $\approx 3.5$ (layer 5). ALiBi exhibits lower
entropy ($\approx 0.3$ less per layer), suggesting its linear bias encourages
sharper positional attention, which may explain its lower translation quality.

\subsection{Learned Gate Values}

On En-De, AdaptiveRoPE gates converge to mean $\approx 0.90$ (query) and
$\approx 0.88$ (key), confirming systematic suppression of high-frequency
components. Per-head analysis reveals specialisation: heads 0 and 1 suppress
most aggressively ($g \approx 0.85$), while head 5 remains closest to 1.0
($g \approx 0.95$), suggesting a long-range vs. short-range division of labour.

%% ─────────────────────────────────────────────────────────────────────────────
\section{Discussion}

Our cross-lingual results challenge the assumption that a single positional
encoding is universally optimal. On high-resource En-De, standard RoPE is
sufficient; learnable adaptations provide marginal gains. But on low-resource
Hi-En, PhasesOnly---which adds only phase offsets without frequency gating---
outperforms RoPE by 1.6 BLEU. This suggests that morphologically rich languages
with flexible word order benefit more from \emph{adjustable positional origins}
than from \emph{frequency scaling}.

The contrast between Hi-En (PhasesOnly best) and Bn-En (AdaptiveRoPE best)
is particularly striking. Both are low-resource Indic languages, yet their
optimal encoding differs. This may reflect linguistic differences: Hindi has
more flexible word order and richer case marking than Bengali, creating
positional patterns that phase offsets can better adapt to.

A limitation is that Hi-En and Bn-En experiments use single seeds due to
compute constraints. Future work should replicate with multiple seeds and
extend to additional language families.

%% ─────────────────────────────────────────────────────────────────────────────
\section{Conclusion}

We presented AdaptiveRoPE and evaluated it across three language pairs. Our
key findings are: (1) the optimal positional encoding is task-dependent;
(2) phase offsets benefit morphologically rich low-resource languages;
(3) frequency gates provide stability on high-resource pairs; and
(4) ALiBi is significantly weaker than RoPE for MT. These results motivate
future work on language-adaptive positional encoding initialisation.

%% ─────────────────────────────────────────────────────────────────────────────
\bibliographystyle{IEEEtran}
\bibliography{references}

\end{document}
